\section{A summable correction for the complete raw source}
\label{sec:v43-complete-correction}
A collision cutoff has no effect on a component whose collision label
is already fixed. We use this elementary fact before extracting the
singularities. It gives one correction for all collision labels, rather
than a sequence of corrections whose germs are chosen anew at every
cutoff. The correction includes the decision and grazing boundaries
through the complete finite graph of Section~\ref{sec:finite-count-extraction}.
Its construction is different from the exponential two-jet word sum
in Section~\ref{sec:v42-morse-packet}.

Fix $n\ge1$ and $R\in I$. Let $w$ be one bounded complex function on
the actual initial section, admissible in the sense of
Definition~\ref{def:finite-record-weight} for every finite cutoff, and
put $M=\|w\|_\infty$. The same function $w$ is used at every cutoff.
Write $\mu^w=(J_{n,R})_*(w\nu_R^*)$. For
$\ell=(k_1,k_2,m)$ set
\begin{equation}\label{eq:v43-component-mass}
 b_\ell=\nu_R^*\{(K_{n,R},N_{n,R})=(k,m)\},\qquad
 f_\ell(t)=\frac{d\mu^w(\ell,\cdot)}{dt}(t).
\end{equation}
Thus $\|f_\ell\|_1\le M b_\ell$ and $\sum_\ell b_\ell=1$.
The full variation convention is used throughout. At a fixed $m$ only
finitely many $k$ occur, and $f_\ell$ is supported in
$[0,m\tau_{\max}]$. Components of zero mass are identically zero.

\subsection{Intrinsic germs with an independently prescribed mass}
\begin{lemma}[Localized singular jets with arbitrary variation budget]
\label{lem:v43-small-jet-mass}
Let $f$ be a compactly supported integrable function with constructible
real and imaginary parts. For every $\beta>0$ one may select a finite
power--logarithm density $e$ such that
\begin{equation}\label{eq:v43-one-component}
 \|e\|_1\le\beta,\qquad q=f-e\in W^{2,1}(\R),\qquad
 \|q\|_1\le\|f\|_1+\beta.
\end{equation}
Every one-sided term of exponent at most one in the intrinsic density
germ is retained. Only its localization radius is changed to meet
$\beta$. In particular this is not an approximation of a singular
coefficient by a smaller coefficient.

The choices can be made once from $f$ and $\beta$, independently of
any collision cutoff. The value and first derivative of $q$ have zero
traces at every extracted point.
\end{lemma}
\begin{proof}
Use the finite analytic partition and the convergent one-sided germs
of Lemma~\ref{lem:power-log-jet}. Equivalently, take the intrinsic
finite set outside which the density has a real-analytic representative,
including the endpoints of its nonzero support. Null values of the
original density do not affect this set or any germ. Denote the jet
at $s$ on side $\sigma$ by
\[
 J_{s,\sigma}(x)=\sum_{\alpha,k}c_{s,\sigma,\alpha,k}
                         x^\alpha(\log x)^k,\qquad -1<\alpha\le1.
\]
These are the germs of the \emph{whole} component. Contributions with
the same roof value have already been added. No division by a
separation between critical values occurs.

Choose one fixed $C^2$ piecewise-polynomial function $\psi$ on
$[0,\infty)$, equal to one on $[0,1/2]$, zero on $[1,\infty)$,
and between zero and one. Its first two derivatives match at the
joining points. For a dyadic $d<1$, the absolute integral of the
localized term is bounded by
\begin{equation}\label{eq:v43-jet-majorant}
 |c_{s,\sigma,\alpha,k}|\int_0^d x^\alpha|\log x|^k\,dx
 =|c_{s,\sigma,\alpha,k}|\,\mathcal J(\alpha+1,k,d).
\end{equation}
Every quantity on the right tends to zero as $d\downarrow0$.
There are finitely many terms. Choose the first dyadic radius small
enough that their total majorant is at most $\beta$, that the
neighborhoods of distinct $s$ are disjoint, and that each neighborhood
lies in its germ interval. One fixes a dyadic subinterval of convergence of each intrinsic
Puiseux--logarithm expansion before this choice. The positive radii
of its coefficient power series, after clearing denominators, determine
such intervals; real analytic continuation alone is not used as a
Taylor convergence radius. None of these choices uses an external cutoff.
For $f=0$ use $e=q=0$.

Set
\[
 e(t)=\sum_{s,\sigma}\one_{\{\sigma(t-s)>0\}}
       \psi\bigl(\sigma(t-s)/d\bigr)J_{s,\sigma}(\sigma(t-s)).
\]
The majorants prove its mass bound. The convergent remainder has
exponent strictly larger than one after the jet is removed. Hence its
value and first derivative vanish at the endpoint and its second
derivative is integrable. On the transition intervals the cutoff is
$C^2$ and the density is analytic. The zero-trace argument of
Proposition~\ref{prop:finite-jet-subtraction} proves $q\in W^{2,1}$.
The final bound is the triangle inequality. Shrinking $d$ can increase
$\|q''\|_1$; no bound for that norm is inferred from $\beta$.
\end{proof}

\subsection{One complete correction, compatible with every count cutoff}
\begin{theorem}[Count-compatible extraction of all boundary germs]
\label{thm:v43-complete-correction}
For every prescribed $0<\varepsilon\le1$ there are finite complex
measures $E^\varepsilon,Q^\varepsilon$ such that
\begin{equation}\label{eq:v43-global-extraction}
 \mu^w=E^\varepsilon+Q^\varepsilon,\qquad
 \|E^\varepsilon\|_{\TV}\le\varepsilon,\qquad
 \|Q^\varepsilon\|_{\TV}\le M+\varepsilon.
\end{equation}
Their densities $e_\ell,q_\ell$ satisfy
\begin{equation}\label{eq:v43-component-budgets}
 \|e_\ell\|_1\le\varepsilon b_\ell,\qquad
 q_\ell\in W^{2,1}(\R),\qquad
 \|q_\ell\|_1\le(M+\varepsilon)b_\ell.
\end{equation}
Each $e_\ell$ is a finite localized power--logarithm sum for the
complete actual component, including singular and section-decision
boundaries. Both series over $\ell$ converge absolutely in total
variation. There is no omitted source in \eqref{eq:v43-global-extraction}.

For every $L\ge n$, define $E^{\varepsilon,L},Q^{\varepsilon,L}$ by
restriction to labels $m\le L$. Then
\begin{equation}\label{eq:v43-projective-compatibility}
 \mu^{w,L}=E^{\varepsilon,L}+Q^{\varepsilon,L},\qquad
 E^{\varepsilon,L'}|_{m\le L}=E^{\varepsilon,L}\quad(L'\ge L),
\end{equation}
and the same restriction identity holds for $Q$. With the cumulative
return constants $a,c_0,C>0$,
\begin{align}
 \|E^\varepsilon-E^{\varepsilon,L}\|_{\TV}
       &\le C\varepsilon e^{an-c_0L},\notag\\
 \|Q^\varepsilon-Q^{\varepsilon,L}\|_{\TV}
       &\le C(M+\varepsilon)e^{an-c_0L}.
                                      \label{eq:v43-correction-count-tail}
\end{align}
For each $0<s<c_0$ there is $C_s<\infty$ such that
\begin{equation}\label{eq:v43-correction-exponential-moment}
 \sum_\ell e^{sm}\|e_\ell\|_1
      \le\varepsilon C_s e^{s(a/c_0)n},\qquad
 \sum_\ell e^{sm}\|q_\ell\|_1
      \le(M+\varepsilon)C_s e^{s(a/c_0)n}.
\end{equation}
The constants in these mass estimates are uniform in $R$.
\end{theorem}
\begin{proof}
For a fixed $\ell=(k,m)$ with $b_\ell>0$, apply the complete finite
collision graph at cutoff $m$ in Lemma~\ref{lem:finite-graph-density}.
The resulting component is exactly $f_\ell$, not a truncated
approximation to it. Indeed the event specifying $N_n=m$ contains
no trajectories with a different collision count. A larger cutoff
therefore gives the identical measure, identical analytic density
pieces and identical intrinsic one-sided germs. This observation is
made before any neighborhood or localization radius is chosen.

Apply Lemma~\ref{lem:v43-small-jet-mass} to $f_\ell$ with
$\beta=\varepsilon b_\ell$. Use zero for a null component. The
component estimates sum to \eqref{eq:v43-global-extraction}.
The measure identity follows component by component; countably many
null sets are harmless. In particular grazing and decision boundaries
are accounted for by their contribution to the finite sublevel
integrals, rather than by a smooth extension of their inverse Jacobian.
The geometric integration input is exactly
\cite[Theorem 1.3 and Theorem 3.11]{CMInt}, as verified in
Section~\ref{sec:finite-count-extraction}.

Use these same component choices for every $L$. This proves
\eqref{eq:v43-projective-compatibility}. Summing
\eqref{eq:v43-component-budgets} over $m>L$ and applying
Theorem~\ref{thm:cumulative-return-tail} proves
\eqref{eq:v43-correction-count-tail}. Applying the exponential-moment
part of that theorem proves \eqref{eq:v43-correction-exponential-moment}.
The construction may vary with $R$, but none of these inequalities
requires continuity of its germ labels or radii in $R$.
\end{proof}

\begin{corollary}[An absolutely convergent correction transform]
\label{cor:v43-correction-transform}
For real frequencies $\omega=(u_1,u_2,s,b)$,
\[
 \widehat E^\varepsilon(\omega)
  =\sum_{\ell=(k,m)}e^{i(u\cdot k+sm)}\widehat e_\ell(b)
\]
converges absolutely and uniformly, and the analogous assertion holds
for $Q^\varepsilon$. The count-truncated transforms have the uniform
errors in \eqref{eq:v43-correction-count-tail}. For each fixed $n,R$
both transforms extend holomorphically to a nonempty complex tube
about the real frequencies. This tube can be chosen uniformly in $R$.
\end{corollary}
\begin{proof}
The first assertion follows from
$|\widehat e_\ell(b)|\le\|e_\ell\|_1$ and the summable majorant.
A nonzero component has $|k|_1\le C_\kappa m$. Its original roof
support is inside $[0,m\tau_{\max}]$, and the localized correction
has support within distance one of that interval. Thus both $e_\ell$
and $q_\ell$ are supported where $|t|\le m\tau_{\max}+1$.
On a compact sub-tube with
\[
 C_\kappa|\operatorname{Im}u|_\infty
 +|\operatorname{Im}s|+\tau_{\max}|\operatorname{Im}b|<c_0,
\]
choose an intermediate $s_0<c_0$ exceeding the left side. The modulus
of each exponential integrand is bounded by a fixed constant times
$e^{s_0m}$. Equation~\eqref{eq:v43-correction-exponential-moment}
and its version with a slightly larger exponent dominate all
fixed-order derivatives as well. Locally uniform convergence of the
entire component transforms proves holomorphy. This is exponential
control in the output variables, not decay as the real roof frequency
tends to infinity.
\end{proof}

\paragraph{What has been summed.}
The series above is a count-labeled series of localized germs of the
whole raw density. It includes fractional powers and logarithms from
boundary sources, and converges without any separation assumption
between moving critical roof values. It is not a proof that the
unlocalized exponential two-jet coefficients of all critical words
are absolutely summable. Nor does \eqref{eq:v43-global-extraction}
bound $\|e^\varepsilon\|_\infty$ or
$\sum_\ell\|q_\ell''\|_1$. In particular it proves a common
all-boundary correction as a finite measure, not the small common
pointwise remainder required by Theorem~\ref{thm:LLT}.
